\documentclass[journal]{IEEEtran}

% packages
\usepackage{ulem}
\usepackage{hyperref}
\usepackage[inline]{enumitem}
\usepackage{array,multirow,graphicx}
\usepackage{float}

%% graphics packages
\usepackage{tikz}
\usepackage{pgfplots}
\usepackage{graphicx}
\usepackage{circuitikz}
\usepackage{dblfloatfix}
\usepackage[bottom]{footmisc}
\ifCLASSOPTIONcompsoc
    \usepackage[caption=false, font=normalsize, labelfont=sf, textfont=sf]{subfig}
\else
    \usepackage[caption=false, font=footnotesize]{subfig}
\fi

%% math packages
\usepackage{amssymb}
\usepackage{amsmath}
\usepackage{mathtools}
\usepackage{booktabs}

% math 
\IEEEeqnarraydefcolsep{0}{\leftmargini}
\newcommand*{\defeq}{\stackrel{\text{def}}{=}}

% correct bad hyphenation here
\hyphenation{op-tical net-works semi-conduc-tor}

% nomenclature
% % Math Operators
% \DeclareMathOperator{\rank}{rank}
% \DeclareMathOperator{\diag}{diag}
% \DeclareMathOperator{\diagm}{diagm}
% \DeclareMathOperator{\real}{Re}
% \DeclareMathOperator{\imag}{Im}
% \DeclareMathOperator{\trace}{tr}
% \DeclareMathOperator{\sign}{sign}
% \DeclareMathOperator{\subjectto}{s.t.}
% \DeclareMathOperator{\matr}{mat}
% \DeclareMathOperator{\vecr}{vec}
% \DeclareMathOperator{\atantwo}{atan2}
% \DeclareMathOperator{\postproc}{post}
% \DeclareMathOperator{\preproc}{pre}
% \DeclareMathOperator{\vars}{var}

% Time
\newcommand{\Time}{t}
\newcommand{\Period}{T}

% Numbers
\newcommand{\NaturalNumberSet}{\mathbb{N}}
\newcommand{\RealNumberSet}{\mathbb{R}}
\newcommand{\ComplexNumberSet}{\mathbb{C}}

% Vector
\newcommand{\im}{j}

% Random Variables and Moments
\newcommand{\RVX}{\mathsf{x}}
\newcommand{\RVY}{\mathsf{y}}
\newcommand{\RVZ}{\mathsf{z}}
\newcommand{\ApproxRVX}{\hat{\RVX}}
\newcommand{\ApproxRVY}{\hat{\RVY}}
\newcommand{\ApproxRVZ}{\hat{\RVZ}}
\newcommand{\Probability}{\mathbb{P}}
\newcommand{\Expectation}{\mathbb{E}}
\newcommand{\Variance}{\mathbb{V}}
\newcommand{\Skewness}{\mathbb{S}}
\newcommand{\Kurtosis}{\mathbb{K}}

% Sets and Indices
\newcommand{\Node}{i}
\newcommand{\AltNode}{j}
\newcommand{\NodeSet}{\mathcal{I}}
\newcommand{\RefNodeSet}{\NodeSet^{\text{ref}}}
\newcommand{\Edge}{e}
\newcommand{\EdgeSet}{\mathcal{E}}
\newcommand{\Line}{b}
\newcommand{\LineSet}{\mathcal{B}}
\newcommand{\Unit}{u}
\newcommand{\UnitSet}{\mathcal{U}}
\newcommand{\PCEIndex}{k}
\newcommand{\PCESet}{\mathcal{K}}
\newcommand{\FixedUnitSet}{\UnitSet^{\text{f}}}
\newcommand{\HarmonicUnitSet}{\UnitSet^{\text{h}}}
\newcommand{\DispatchableUnitSet}{\UnitSet^{\text{d}}}
\newcommand{\Harmonic}{h}
\newcommand{\HarmonicSet}{\mathcal{H}}
\newcommand{\PositiveHarmonicSet}{\HarmonicSet^{\texttt{+}}}
\newcommand{\NegativeHarmonicSet}{\HarmonicSet^{\texttt{-}}}
\newcommand{\ZeroHarmonicSet}{\HarmonicSet^{\texttt{0}}}
\newcommand{\Winding}{w}
\newcommand{\WindingSet}{\mathcal{W}}
\newcommand{\Xfmr}{x}
\newcommand{\XfmrSet}{\mathcal{X}}


% Union Sets and Indices
\newcommand{\TupleSet}{\mathcal{T}}
\newcommand{\LineTuple}{\Line\Node\AltNode}
\newcommand{\AltLineTuple}{\Line\AltNode\Node}
\newcommand{\EdgeTuple}{\Edge\Node\AltNode}
\newcommand{\AltEdgeTuple}{\Edge\AltNode\Node}
\newcommand{\EdgeTupleSet}{\TupleSet^{\text{\Edge}}}
\newcommand{\EdgeTupleSetFr}{\TupleSet^{\text{\Edge}\Arrow{.15cm}}}
\newcommand{\EdgeTupleSetTo}{\TupleSet^{\text{\Edge}\Arrowright{.15cm}}}
\newcommand{\LineTupleSet}{\TupleSet^{\text{\Line}}}
\newcommand{\LineTupleSetFr}{\TupleSet^{\text{\Line}\Arrow{.15cm}}}
\newcommand{\LineTupleSetTo}{\TupleSet^{\text{\Line}\Arrowright{.15cm}}}
\newcommand{\XfmrTupleSet}{\TupleSet^{\text{\Xfmr}}}
\newcommand{\XfmrTupleSetFr}{\TupleSet^{\text{\Xfmr}\Arrow{.15cm}}}
\newcommand{\XfmrTupleSetTo}{\TupleSet^{\text{\Xfmr}\Arrowright{.15cm}}}
\newcommand{\UnitTuple}{\Unit\Node}
\newcommand{\UnitTupleSet}{\TupleSet^{\text{\Unit}}}
\newcommand{\FixedUnitTupleSet}{\TupleSet^{\text{\Unit,f}}}
\newcommand{\HarmonicUnitTupleSet}{\TupleSet^{\text{\Unit,h}}}
\newcommand{\DispatchableUnitTupleSet}{\TupleSet^{\text{\Unit,d}}}

\newcommand{\Arrow}[1]{%
\parbox{#1}{\tikz{\draw[->](0,0)--(#1,0);}}
}
\newcommand{\Arrowright}[1]{%
\parbox{#1}{\tikz{\draw[<-](0,0)--(#1,0);}}
}

% Electrical Variables
%% Current
\newcommand{\Current}{\mathsf{I}}
\newcommand{\ComplexCurrent}{\mathsf{\bar{I}}}
\newcommand{\CurrentMagnitude}{\Current^{\text{mag}}}
\newcommand{\RealCurrent}{\Current^{\text{re}}}
\newcommand{\ImaginaryCurrent}{\Current^{\text{im}}}
\newcommand{\SeriesCurrent}{\Current^{\text{s}}}
\newcommand{\ComplexSeriesCurrent}{\ComplexCurrent^{\text{s}}}
\newcommand{\RealSeriesCurrent}{\Current^{\text{s,re}}}
\newcommand{\RealSeriesCurrentFrom}{\Current^{\text{s,re,fr}}}
\newcommand{\RealSeriesCurrentTo}{\Current^{\text{s,re,to}}}
\newcommand{\ImaginarySeriesCurrent}{\Current^{\text{s,im}}}
\newcommand{\ImaginarySeriesCurrentFrom}{\Current^{\text{s,im,fr}}}
\newcommand{\ImaginarySeriesCurrentTo}{\Current^{\text{s,im,to}}}
\newcommand{\ShuntCurrent}{\Current^{\text{sh}}}
\newcommand{\ComplexShuntCurrent}{\ComplexCurrent^{\text{sh}}}
\newcommand{\RealShuntCurrent}{\Current^{\text{sh,re}}}
\newcommand{\ImaginaryShuntCurrent}{\Current^{\text{sh,im}}}
\newcommand{\RealMagnetizingCurrent}{\Current^{\text{m,re}}}
\newcommand{\ImaginaryMagnetizingCurrent}{\Current^{\text{m,im}}}
\newcommand{\RealExcitationCurrent}{\Current^{\text{e,re}}}
\newcommand{\ImaginaryExcitationCurrent}{\Current^{\text{e,im}}}

\newcommand{\CurrentWave}[1]{i_{#1}(\Time)}
%% Voltage
\newcommand{\Voltage}{\mathsf{U}}
\newcommand{\ComplexVoltage}{\mathsf{\bar{U}}}
\newcommand{\RealVoltage}{\Voltage^{\text{re}}}
\newcommand{\ImaginaryVoltage}{\Voltage^{\text{im}}}
\newcommand{\ReferenceVoltage}{\Voltage^{\text{ref}}}

\newcommand{\AltVoltageVariable}{\mathsf{V}}
\newcommand{\ComplexAltVoltage}{\mathsf{\bar{V}}}
\newcommand{\RealAltVoltage}{\AltVoltageVariable^{\text{re}}}
\newcommand{\ImaginaryAltVoltage}{\AltVoltageVariable^{\text{im}}}

\newcommand{\LiftedCurrentVariable}{\mathsf{J}}
\newcommand{\LiftedCurrent}{\LiftedCurrentVariable}
\newcommand{\LiftedVoltageVariable}{\mathsf{W}}
\newcommand{\LiftedVoltage}{\LiftedVoltageVariable}

\newcommand{\VoltageWave}[1]{u_{#1}(\Time)}

%% Excitation Voltage
\newcommand{\ExcitationVoltageVariable}{\mathsf{E}}
\newcommand{\ComplexExcitationVoltage}{\mathsf{\bar{E}}}
\newcommand{\RealExcitationVoltage}{\ExcitationVoltageVariable^{\text{re}}}
\newcommand{\ImaginaryExcitationVoltage}{\ExcitationVoltageVariable^{\text{im}}}
\newcommand{\RealExcitationVoltageVector}{\overline{\ExcitationVoltageVariable}^{\text{re}}}
\newcommand{\ImaginaryExcitationVoltageVector}{\overline{\ExcitationVoltageVariable}^{\text{im}}}

%% Power 
\newcommand{\Power}{\mathsf{S}}
\newcommand{\RealPower}{\mathsf{P}}
\newcommand{\ImaginaryPower}{\mathsf{Q}}

% Electrical Parameters
%% Frequency
\newcommand{\Frequency}{f}
\newcommand{\FundamentalFrequency}{\Frequency^{\text{fund}}}
\newcommand{\AngularFrequency}{\omega}
\newcommand{\FundamentalAngularFrequency}{\AngularFrequency^{\text{fund}}}
%% Current
\newcommand{\CurrentParameter}{I}
\newcommand{\CurrentPhasor}[1]{\CurrentParameter_{#1}}
\newcommand{\SeriesCurrentPhasor}[1]{\CurrentPhasor{#1}^{\text{s}}}
\newcommand{\ShuntCurrentPhasor}[1]{\CurrentPhasor{#1}^{\text{sh}}}
\newcommand{\ExcitationCurrentPhasor}[1]{\CurrentParameter_{#1}^{\text{e}}}
\newcommand{\MagnetizingCurrentPhasor}[1]{\CurrentParameter_{#1}^{\text{m}}}
\newcommand{\CurrentMagnitude}[1]{\CurrentParameter_{#1}^{\text{mag}}}
\newcommand{\RatedCurrent}{\CurrentParameter^{\text{rated}}}
\newcommand{\FundamentalCurrentMagnitude}[1]{\CurrentParameter_{#1}^{\text{mag,fund}}}

\newcommand{\CurrentAngleParameter}{\varphi}
\newcommand{\CurrentAngle}[1]{\CurrentAngleParameter_{#1}}

%% Voltage
\newcommand{\VoltageParameter}{U}
\newcommand{\MinimumVoltage}{\VoltageParameter^{\text{min}}}
\newcommand{\MaximumVoltage}{\VoltageParameter^{\text{max}}}
\newcommand{\ExtremeVoltage}{\hat{\VoltageParameter}^{\text{max}}}
\newcommand{\VoltageMagnitude}[1]{\VoltageParameter_{#1}^{\text{mag}}}
\newcommand{\FundamentalVoltageMagnitude}[1]{\VoltageParameter_{#1}^{\text{mag,fund}}}
\newcommand{\RealRefVoltage}{\VoltageParameter^{\text{ref,re}}}
\newcommand{\ImaginaryRefVoltage}{\VoltageParameter^{\text{ref,im}}}

\newcommand{\AltVoltageParameter}{V}

\newcommand{\ExcitationVoltageParameter}{E}
\newcommand{\ExcitationVoltageVector}{\overline{\ExcitationVoltageParameter}}
\newcommand{\ExcitationVoltagePhasor}[1]{\ExcitationVoltageParameter_{#1}}

\newcommand{\VoltageAngleParameter}{\theta}
\newcommand{\VoltageAngle}[1]{\VoltageAngleParameter_{#1}}

\newcommand{\VoltagePhasor}[1]{\VoltageParameter_{#1}}
\newcommand{\AltVoltagePhasor}[1]{\AltVoltageParameter_{#1}}
%% Power
\newcommand{\ApparentPowerParameter}{S}
\newcommand{\ApparentPowerPhasor}[1]{\ApparentPowerParameter_{#1}}

\newcommand{\RealPowerParameter}{P}
\newcommand{\MinimumRealPower}{P^{\text{min}}}
\newcommand{\MaximumRealPower}{P^{\text{max}}}
\newcommand{\ReferenceRealPower}{P^{\text{ref}}}
\newcommand{\ImaginaryPowerParameter}{Q}
\newcommand{\MinimumImaginaryPower}{Q^{\text{min}}}
\newcommand{\MaximumImaginaryPower}{Q^{\text{max}}}
\newcommand{\ReferenceImaginaryPower}{Q^{\text{ref}}}
\newcommand{\PowerFactor}{\theta}

%% Impedance
\newcommand{\Impedance}{\bar{z}}
\newcommand{\SeriesImpedance}[1]{\Impedance_{#1}^{\text{s}}}
\newcommand{\Resistance}{r}
\newcommand{\SeriesResistance}[1]{\Resistance_{#1}^{\text{s}}}
\newcommand{\ShuntResistance}[1]{\Resistance_{#1}^{\text{sh}}}
\newcommand{\Capacitance}{C}
\newcommand{\SeriesCapacitance}[1]{\Capacitance_{#1}^{\text{s}}}
\newcommand{\Inductance}{L}
\newcommand{\SeriesInductance}[1]{\Inductance_{#1}^{\text{s}}}
\newcommand{\Reactance}{x}
\newcommand{\SeriesReactance}[1]{\Reactance_{#1}^{\text{s}}}
\newcommand{\Admittance}{\bar{y}}
\newcommand{\ShuntAdmittance}[1]{\Admittance_{#1}^{\text{sh}}}
\newcommand{\Conductance}{g}
\newcommand{\ShuntConductance}[1]{\Conductance_{#1}^{\text{sh}}}
%\newcommand{\ShuntConductance}{\Conductance^{\text{sh}}}
\newcommand{\Susceptance}{b}
\newcommand{\ShuntSusceptance}[1]{\Susceptance_{#1}^{\text{sh}}}
%\newcommand{\ShuntSusceptance}{\Susceptance^{\text{sh}}}

%% Transformation
\newcommand{\Transform}{t}
\newcommand{\RealTransform}{\Transform^{\text{re}}}
\newcommand{\ImaginaryTransform}{\Transform^{\text{im}}}
\newcommand{\VectorGroupTransform}{\Transform^{\text{vg}}}
\newcommand{\RealVectorGroupTransform}{\Transform^{\text{vg,re}}}
\newcommand{\ImaginaryVectorGroupTransform}{\Transform^{\text{vg,im}}}

%% Harmonic injection limits
\newcommand{\Hlimit}{m}
\newcommand{\RealHlimit}{m^{\text{re}}}
\newcommand{\ImaginaryHlimit}{m^{\text{im}}}
%% wye
\newcommand{\wye}{\mathbin{\tikz[x=1ex,y=1ex]{\draw[line width=.1ex] (0,0)--(45:1)--++(-45:1) (45:1)--++(0,1);}}}

% Polynomial Chaos Expansion
\newcommand{\CoefficientX}{x}
\newcommand{\CoefficientY}{y}
\newcommand{\CoefficientZ}{z}
\newcommand{\Polynoom}{\psi}

%# Inner product
\newcommand{\SecondIP}{\langle \Polynoom_{\Node}, \Polynoom_{\Node} \rangle}
\newcommand{\ThirdIP}{\langle \Polynoom_{\Node_{1}}, \Polynoom_{\Node_{2}}, \Polynoom_{\Node} \rangle}

\newcommand{\psione}{ \Polynoom_{\Node_{1}}}
\newcommand{\psitwo}{ \Polynoom_{\Node_{2}}}

\newcommand{\norm}[1]{\left\lVert#1\right\rVert}
\newcommand*{\gp}{\stackrel{\text{gp}}{\to}}
\newcommand*{\downgp}{\downarrow\text{\footnotesize gp}}

% document
\begin{document}

\section{Mathematical Framework}
\label{sec:mathematical_framework}

This section describes the real-value feasible set for the harmonic hosting capacity in the rectangular current-voltage formulation of the harmonic optimal power flow problem. The following assumptions are made:
\begin{enumerate}
    \item a balanced three-phase power system is considered, i.e., the harmonics are classified into three types:
    \begin{enumerate}
        \item positive sequence~$\Harmonic \in \PositiveHarmonicSet \subset \HarmonicSet$:~$\Harmonic=1+3n$,
        \item negative sequence~$\Harmonic \in \NegativeHarmonicSet \subset \HarmonicSet$:~$\Harmonic=2+3n$, and
        \item zero sequence~$\Harmonic \in \ZeroHarmonicSet \subset \HarmonicSet$ :~$\Harmonic=3+3n$,
    \end{enumerate}
    where~$n \in \NaturalNumberSet_{0}$;
    \item any edge in the system is either represented as a pi-section or a t-section, with parameters~$r$, $x$, $g$ and $b$;
    \item the fundamental power planning problem is out of scope, consequently the fundamental harmonic is omitted throughout the presentation, as such: $\HarmonicSet \defeq \HarmonicSet \backslash \{1\}$;
    \item the fundamental voltage magnitude at a node~$\Node \in \NodeSet$ becomes a deterministic parameter $\VoltageMagnitude{\Node,1}$ and is set to 1.0\,pu;
    \item the lower bound on the non-fundamental voltage magnitude at any node is set to zero, consequently, the lower bound on the root mean square voltage expires; 
    \item the lower bound on the non-fundamental generator active and reactive power is set to zero; and
    % \item any unit in the system is represented as a current source, more complex unit models can be added;
    % \item the fundamental voltage magnitude at any bus~$|U|_{i,1},~\forall i \in I$ is set to 1\,p.u.;
    \item the rectangular non-fundamental harmonic voltage components of the reference bus(ses)~$U^{re}_{i,h},~U^{im}_{i,h},~\forall i \in \RefNodeSet, h \in \HarmonicSet \backslash \{1\}$ are set to 0\,p.u.
\end{enumerate}

% Note that for the root mean square voltage, the customary lower bound is omitted as it becomes redundant given assumption four. 

% The following assumptions are made:
% \begin{enumerate}
%     \item The voltage magnitude of the fundamental harmonic, i.e., $\Harmonic = 1$, is treated as an uncertain variable~$\LiftedVoltage_{\Node,1}$. Consequently, the right hand side of \eqref{eq:thd} and \eqref{eq:ihd} is simplified.
% \end{enumerate}

Following abbreviations are used: $\alpha^{\text{x,y}} = y(x(\alpha^{\text{min}}),x(\alpha^{\text{max}}))$, e.g., $\alpha^{\text{cos,min}} = \min(\sin(\alpha^{\text{min}},\alpha^{\text{max}}))$

{\allowdisplaybreaks 
\begin{IEEEeqnarray}{ 0 l }
    \textsc{objective function} \nonumber \\
    \quad \max \sum_{\Unit \in \FixedUnitSet}\sum_{\Harmonic \in \HarmonicSet} \CurrentMagnitude_{\Unit,\Harmonic} \\
    \textsc{Bus Constraints} \nonumber \\
    \text{reference bus --~} \forall \Node \in \RefNodeSet: \nonumber\\
    \quad \RealVoltage_{\Node,\Harmonic} = 0, \,\, \forall\Harmonic \in \HarmonicSet \label{eq_ref_bus_volt_source_a}, \\ 
    \quad \ImaginaryVoltage_{\Node,\Harmonic} = 0, \,\, \forall \Harmonic \in \HarmonicSet, \label{eq_ref_bus_volt_source_b} \\
    \text{Kirchhoff's current law --~} \forall \Node \in \NodeSet, \Harmonic \in \HarmonicSet: \nonumber \\
    \quad \sum_{\LineTuple \in \LineTupleSet} \mkern-6mu \RealCurrent_{\LineTuple,\Harmonic} + \mkern-9mu
    \sum_{\UnitTuple \in \UnitTupleSet} \mkern-6mu \RealCurrent_{\UnitTuple,\Harmonic}  + \ShuntConductance{\Node,\Harmonic} \RealVoltage_{\Node,\Harmonic} -\ShuntSusceptance{\Node,\Harmonic} \ImaginaryVoltage_{\Node,\Harmonic}  = 0, \label{eq:kcl_node_real} \\
    \quad \sum_{\LineTuple \in \LineTupleSet} \mkern-6mu \ImaginaryCurrent_{\LineTuple,\Harmonic} + \mkern-9mu
    \sum_{\UnitTuple \in \UnitTupleSet} \mkern-6mu \ImaginaryCurrent_{\UnitTuple,\Harmonic} + \ShuntConductance{\Node,\Harmonic} \ImaginaryVoltage_{\Node,\Harmonic} + \ShuntSusceptance{\Node,\Harmonic} \RealVoltage_{\Node,\Harmonic}= 0, \label{eq:kcl_node_imag} \\
    \text{root mean square voltage limit --~} \forall \Node \in \NodeSet: \nonumber \\
    \quad \sum_{\Harmonic \in \HarmonicSet} (\RealVoltage_{\Node,\Harmonic})^{2} + (\ImaginaryVoltage_{\Node,\Harmonic})^{2} \leq (\MaximumVoltage_{\Node})^{2} - (\VoltageMagnitude{\Node,1})^{2}, \label{eq:voltage_bounds_fund} \\
    \text{total harmonic voltage limit --~} \forall \Node \in \NodeSet: \nonumber \\
    \quad \sum_{\Harmonic \in \HarmonicSet} (\RealVoltage_{\Node,\Harmonic})^{2} + (\ImaginaryVoltage_{\Node,\Harmonic})^{2} \leq (\VoltageMagnitude{\Node,1} \, \text{THD}_{\Node}^{\text{u}})^2, \label{eq:thd_bounds} \\
    \text{individual harmonic voltage limit --~} \forall \Node \in \NodeSet, \Harmonic \in \HarmonicSet: \nonumber \\
    \quad (\RealVoltage_{\Node,\Harmonic})^{2} + (\ImaginaryVoltage_{\Node,\Harmonic})^{2} \leq (\VoltageMagnitude{\Node,1} \, \text{IDH}^{\text{u}}_{\Node,\Harmonic})^{2},
    \label{eq:harmonic_bounds} \\
    \textsc{Branch Constraints} \nonumber \\
    \text{Ohm's law --~} \forall \LineTuple \in \LineTupleSetFr, \Harmonic \in \HarmonicSet: \nonumber\\
    \quad \RealVoltage_{\AltNode,\Harmonic} = \RealVoltage_{\Node,\Harmonic} - 
    \SeriesResistance{\Line,\Harmonic} \RealSeriesCurrent_{\LineTuple,\Harmonic} +
    \SeriesReactance{\Line,\Harmonic} \ImaginarySeriesCurrent_{\LineTuple,\Harmonic}, \label{eq:ol_real} \\ 
    \quad \ImaginaryVoltage_{\AltNode,\Harmonic} = \ImaginaryVoltage_{\Node,\Harmonic} - 
    \SeriesResistance{\Line,\Harmonic} \ImaginarySeriesCurrent_{\LineTuple,\Harmonic} -
    \SeriesReactance{\Line,\Harmonic} \RealSeriesCurrent_{\LineTuple,\Harmonic}, \label{eq:ol_imag} \\
    \text{branch total current --~} \forall \LineTuple \in \LineTupleSet, \Harmonic \in \HarmonicSet: \nonumber \\
    \quad \RealCurrent_{\LineTuple,\Harmonic} =
    \ShuntConductance{\LineTuple,\Harmonic} \RealVoltage_{\Node,\Harmonic} - 
    \ShuntSusceptance{\LineTuple,\Harmonic} \ImaginaryVoltage_{\Node,\Harmonic} +
    \RealSeriesCurrent_{\LineTuple,\Harmonic}, \label{eq:btc_real} \\
    \quad \ImaginaryCurrent_{\LineTuple,\Harmonic} =
    \ShuntConductance{\LineTuple,\Harmonic} \ImaginaryVoltage_{\Node,\Harmonic} + 
    \ShuntSusceptance{\LineTuple,\Harmonic} \RealVoltage_{\Node,\Harmonic} +
    \ImaginarySeriesCurrent_{\LineTuple,\Harmonic}, \label{eq:btc_imag} \\
    \text{branch current limits --~} \forall \LineTuple \in \LineTupleSet: \nonumber \\
    \quad \sum_{\Harmonic \in \HarmonicSet} (\RealCurrent_{\LineTuple,\Harmonic})^{2} + (\ImaginaryCurrent_{\LineTuple,\Harmonic})^{2} \leq (\RatedCurrent_{\LineTuple})^{2}, \\
    \textsc{Transformer Constraints} \nonumber \\
    \text{xfmr core --~} \forall \Xfmr \in \XfmrSet, \Harmonic \in \HarmonicSet: \nonumber \\
    \quad \RealMagnetizingCurrent_{\Xfmr,\Harmonic} +  \frac{\RealExcitationVoltage_{\Xfmr,\Harmonic}}{\ShuntResistance{\Xfmr,\Harmonic}} =
    \sum_{\Winding \in \WindingSet} \RealVectorGroupTransform_{\Xfmr,\Harmonic}\,\RealSeriesCurrent_{\Winding,\Harmonic} + \ImaginaryVectorGroupTransform_{\Xfmr,\Harmonic}\,\ImaginarySeriesCurrent_{\Winding,\Harmonic}, \\
    \quad \ImaginaryMagnetizingCurrent_{\Xfmr,\Harmonic} +  \frac{\ImaginaryExcitationVoltage_{\Xfmr,\Harmonic}}{\ShuntResistance{\Xfmr,\Harmonic}} =
    \sum_{\Winding \in \WindingSet} \RealVectorGroupTransform_{\Xfmr,\Harmonic}\,\ImaginarySeriesCurrent_{\Winding,\Harmonic} - \ImaginaryVectorGroupTransform_{\Xfmr,\Harmonic}\,\RealSeriesCurrent_{\Winding,\Harmonic}, \\
    \quad \RealAltVoltage_{\Xfmr,\Harmonic,1} = \RealExcitationVoltage_{\Xfmr,\Harmonic} - \SeriesReactance{\Xfmr,\Harmonic} \,\ImaginarySeriesCurrent_{\Xfmr,\Harmonic,1},\\
    \quad \ImaginaryAltVoltage_{\Xfmr,\Harmonic,1} = \ImaginaryExcitationVoltage_{\Xfmr,\Harmonic} + \SeriesReactance{\Xfmr,\Harmonic}\,\RealSeriesCurrent_{\Xfmr,\Harmonic,1}, \\
    \quad \RealAltVoltage_{\Xfmr,\Harmonic,2} = \RealVectorGroupTransform_{\Xfmr,\Harmonic}\,\RealExcitationVoltage_{\Xfmr,\Harmonic} - \ImaginaryVectorGroupTransform_{\Xfmr,\Harmonic}\,\ImaginaryExcitationVoltage_{\Xfmr,\Harmonic}, \\
    \quad \ImaginaryAltVoltage_{\Xfmr,\Harmonic,2} = \RealVectorGroupTransform_{\Xfmr,\Harmonic}\,\ImaginaryExcitationVoltage_{\Xfmr,\Harmonic} + \ImaginaryVectorGroupTransform_{\Xfmr,\Harmonic}\,\RealExcitationVoltage_{\Xfmr,\Harmonic}, \\
    \text{xfmr winding  --~} \forall \Xfmr \in \XfmrSet, \Winding \in \WindingSet_{\Xfmr}, \Harmonic \in \HarmonicSet: \nonumber \\
    \quad \RealSeriesCurrent_{\Winding,\Harmonic} = \RealCurrent_{\Winding,\Harmonic} + 
    \ShuntConductance{\Winding,\Harmonic}\,\RealAltVoltage_{\Winding,\Harmonic} - 
    \ShuntSusceptance{\Winding,\Harmonic}\,\ImaginaryAltVoltage_{\Winding,\Harmonic}, \\
    \quad \ImaginarySeriesCurrent_{\Winding,\Harmonic} = \ImaginaryCurrent_{\Winding,\Harmonic} + \ShuntConductance{\Winding,\Harmonic}\,\ImaginaryAltVoltage_{\Winding,\Harmonic} + 
    \ShuntSusceptance{\Winding,\Harmonic}\,\RealAltVoltage_{\Winding,\Harmonic}, \\
    \text{xfmr winding  --~} \forall \Xfmr \in \XfmrSet, \Winding \in \WindingSet_{\Xfmr}, \Harmonic \in \PositiveHarmonicSet \cup \NegativeHarmonicSet: \Winding \to \Node: \nonumber\\
    \quad \RealVoltage_{\Node,\Harmonic}  = \RealAltVoltage_{\Winding,\Harmonic} + \Resistance_{\Winding,\Harmonic}\,\RealCurrent_{\Winding,\Harmonic},  \\
    \quad \ImaginaryVoltage_{\Node,\Harmonic} = \ImaginaryAltVoltage_{\Winding,\Harmonic} + \Resistance_{\Winding,\Harmonic}\,\ImaginaryCurrent_{\Winding,\Harmonic}, \\ 
    \text{xfmr winding  --~} \forall \Xfmr \in \XfmrSet, \Winding \in \WindingSet^{\text{e}}_{\Xfmr}, \Harmonic \in \ZeroHarmonicSet: \Winding \to \Node: \nonumber \\
    \quad \RealVoltage_{\Node,\Harmonic}  = \RealAltVoltage_{\Winding,\Harmonic} + (r_{\Winding,\Harmonic}+3r^{\text{0}}_{\Winding,\Harmonic})\,\RealCurrent_{\Winding,\Harmonic} - 3x^{\text{0}}_{\Winding,\Harmonic}\,\ImaginaryCurrent_{\Winding,\Harmonic},  \nonumber \\ \\
    \quad \ImaginaryVoltage_{\Node,\Harmonic} = \ImaginaryAltVoltage_{\Winding,\Harmonic} + (r_{\Winding,\Harmonic}+3r^{\text{0}}_{\Winding,\Harmonic})\,\ImaginaryCurrent_{\Winding,\Harmonic} + 3x^{\text{0}}_{\Winding,\Harmonic}\,\RealCurrent_{\Winding,\Harmonic}, \nonumber \\ \\
    \text{xfmr winding  --~} \forall \Xfmr \in \XfmrSet, \Winding \in \WindingSet^{\text{ne}}_{\Xfmr} \cup \WindingSet^{\text{d}}_{\Xfmr}, \Harmonic \in \ZeroHarmonicSet: \Winding \to \Node: \nonumber \\
    \quad \RealCurrent_{\Winding,\Harmonic} = 0, \\ 
    \quad \ImaginaryCurrent_{\Winding,\Harmonic} = 0, \\
    \label{eq_zero_se_block}
    \textsc{Harmonic Load Constraints} \nonumber \\
    \text{load current angle cone --~} \forall \UnitTuple \in \FixedUnitTupleSet, \Harmonic \in \HarmonicSet: \nonumber \\
    \quad \alpha^{\text{cos,min}} \, \CurrentMagnitude_{\Unit,\Harmonic} \leq 
    \RealCurrent_{\Unit,\Harmonic} \leq 
    \alpha^{\text{cos,max}} \, \CurrentMagnitude_{\Unit,\Harmonic} \\
    \quad \alpha^{\text{sin,min}} \, \CurrentMagnitude_{\Unit,\Harmonic} \leq 
    \ImaginaryCurrent_{\Unit,\Harmonic} \leq 
    \alpha^{\text{sin,max}} \, \CurrentMagnitude_{\Unit,\Harmonic} \\
    \quad (\CurrentMagnitude_{\Unit,\Harmonic})^{2} \leq 
    (\RealCurrent_{\Unit,\Harmonic})^2 + (\ImaginaryCurrent_{\Unit,\Harmonic})^2 
\end{IEEEeqnarray}}

\end{document}


